\documentclass{article}
\usepackage[T1]{fontenc}

\title{Assignment 2}

\begin{document}
\maketitle

\section{Part 1}
For a particle of mass $m$ moving with a velocity $v$, the action $S$ along a path is given by the integral of the Lagrangian. Since the potential energy here is $0$, the kinetic energy is given by:
\[
  \frac{1}{2}mv^2
\]
The Lagrangian is:
\[
  L = T - V = \frac{1}{2}mv^2
\]
The action is:
\[
  S = \int L dt = \int \frac{1}{2}mv^2 dt = \frac{1}{2}mv^2 \left(\frac{s}{v}\right) = \frac{1}{2}mvs
\]
Path B is slightly longer than Path A because the neutron travels an extra distance $d$. in Figure 1, that extra distance forms a right triangle with the nucleus's diameter $D$ and the scattering angle $\theta$, giving 
\[
d = D \sin\theta
\]
Because action depends directly on distance, the difference in action between the two paths is:
\[
\Delta S = mvD \sin\theta
\]

\section{Part 2}

Total wave probability amplitude $W_{A\rightarrow B}$ is the sum of the wave functions for path A and path B
    \[
    W_{A\rightarrow B} = W_A + W_B
    \]
    
The probability amplitude for a particle traveling along a specific classical path is:
    \[
    W_A = N e^{iS_A/\hbar}, \quad w_B = N e^{iS_B/\hbar}
    \]
    where $N$ is a normalization constant.
    
The probability $P$ is given by the squared absolute value of the total amplitude:
    \[
    P_{A\rightarrow B} = \left( N e^{iS_A/\hbar} + N e^{iS_B/\hbar} \right)^* \left( N e^{iS_A/\hbar} + N e^{iS_B/\hbar} \right)
    \]
    \[
    P_{A\rightarrow B} = N^2 \left( e^{-iS_A/\hbar} + e^{-iS_B/\hbar} \right) \left( e^{iS_A/\hbar} + e^{iS_B/\hbar} \right)
    \]
    \[
    P{A\rightarrow B} = N^2 \left( 1 + e^{i(S_B - S_A)/\hbar} + e^{-i(S_B - S_A)/\hbar} + 1 \right)
    \]
    but $\Delta S = S_B - S_A$:
    \[
    P{A\rightarrow B} = N^2 \left( 2 + e^{i\Delta S/\hbar} + e^{-i\Delta S/\hbar} \right)
    \]
    
Applying the identity $\cos z = \frac{e^{iz} + e^{-iz}}{2}$:  
    \[
    P = 2N^2 \left( 1 + \cos\frac{\Delta S}{\hbar} \right)
    \]
    Thus the probability to detect a neutron is proportional to $(1 + \cos \frac{\Delta S}{\hbar})$

\section {Part 3}

The detection probability depends on the term $(1 + \cos \frac{\Delta S}{\hbar})$. A destructive interference minimum occurs when the cosine term reaches its minimum value of $-1$. This requires the argument of cosine:
    \[
    \frac{\Delta S}{\hbar} = \pi \Rightarrow |\Delta S| = \pi\hbar
    \]
At this angle ($\theta \approx 18^\circ$), $\Delta S = \pi\hbar$, meaning the phase difference between the two path amplitudes is $\pi$ radians ($180^\circ$). Also, the two phasors associated with paths $A$ and $B$ point in opposite directions, leading to destructive interference. 

\section {Part 4}

At the first minimum angle $\theta_0$, the action difference satisfies $\Delta S = \pi\hbar$[cite: 1]. Substituting for $\Delta S$:
    \[
    m v D \sin\theta_0 = \pi\hbar
    \]
    Solving for $D$:
    \[
    D = \frac{\pi\hbar}{m v \sin\theta_0}
    \]
    Since $\hbar = \frac{h}{2\pi}$, substituting this gives:
    \[
    D = \frac{h}{2 m v \sin\theta_0}
    \]
Given:

Planck constant $h \approx 6.63 \times 10^{-34} J\cdot{s}$

neutron mass $m = 1.67 \times 10^{-27} { kg}$

neutron velocity $v = 10^8 { m/s}$

angle $\theta_0 = 18^\circ$, $\sin(18^\circ) \approx 0.309017$:
    \[
    D = \frac{6.63 \times 10^{-34}}{2 \times (1.67 \times 10^{-27}) \times (10^8) \times 0.309017} = \frac{6.63 \times 10^{-34}}{1.03218 \times 10^{-19}}
    \]
    \[
    D \approx 6.42 \times 10^{-15} { m}
    \]



\end{document}